\documentclass{article}

% For arXiv preprint:
\usepackage[preprint]{neurips_2024}

\usepackage[utf8]{inputenc}
\usepackage[T1]{fontenc}
\usepackage{hyperref}
\usepackage{url}
\usepackage{booktabs}
\usepackage{amsfonts}
\usepackage{amsmath}
\usepackage{nicefrac}
\usepackage{microtype}
\usepackage{xcolor}
\usepackage{enumitem}

\title{Memory-Gravity Attention: Persistent-Importance Retrieval\\for Long-Horizon Memory Tasks}

\author{
  Kevin T.N \\
  Independent Researcher \\
  \texttt{github.com/jkdkr2439/Memory-Gravity-Attention}
}

\begin{document}

\maketitle

\begin{abstract}
We present Memory-Gravity Attention (MGA), a model-agnostic retrieval layer that augments similarity-based memory retrieval with persistent-importance signals. Standard approaches---including similarity-only and Generative-Agents-style scoring (recency + importance + similarity)---perform well on recency-friendly queries but fail on tasks requiring long-horizon persistent memory: retrieving facts stated once long ago (needle), tracking unresolved tasks (open-loop), and filtering noise. MGA introduces a multiplicative gate: $\text{score} = \text{sim} \times (1 + \sigma(\boldsymbol{\theta}^\top \mathbf{f}))$, where $\mathbf{f}$ comprises log-estimated persistent signals and $\boldsymbol{\theta}$ is learned via logistic regression. On a synthetic benchmark of 10 worlds (4{,}500+ nodes, 415 test queries, sentence-transformer embeddings), the Generative-Agents baseline achieves Recall@5 = 0.209 overall but scores 0.000 on needle, open-loop, and noise-resistance query families. MGA achieves non-zero recall on all three (needle: 0.143, open-loop: 0.222, noise: 0.075) while maintaining the lowest noise retrieval rate (Noise@5 = 0.017). An attention-entropy diagnostic confirms no pathological lock behavior. Code available at \url{https://github.com/jkdkr2439/Memory-Gravity-Attention}.
\end{abstract}

\section{Introduction}

Memory retrieval for language model agents is typically framed as a similarity search problem: given a query, retrieve the most semantically similar stored items. This works well when the answer is recent and phrased similarly to the query. However, it fails systematically on tasks requiring \emph{persistent importance}---information that remains relevant over long horizons regardless of recency or surface similarity.

Consider the following scenarios:
\begin{itemize}[leftmargin=2em,itemsep=2pt]
\item A fact stated once in session 2 of 25, never repeated, but still needed (\textbf{needle}).
\item A task opened early and never closed, requiring recall despite intervening events (\textbf{open-loop}).
\item A constraint revised mid-project, where the old version has high similarity to the query but is stale (\textbf{stale trap}).
\item A dense log with 40\%+ chatter, where the system must retrieve signal from noise (\textbf{noise resistance}).
\end{itemize}

Park et al.~\citep{park2023generative} proposed scoring memory as $\alpha \cdot \text{recency} + \beta \cdot \text{importance} + \gamma \cdot \text{relevance}$ for generative agents. This addresses some failure modes but, as we show empirically, still scores zero on needle, open-loop, and noise-resistance tasks.

We propose \textbf{Memory-Gravity Attention (MGA)}, which uses a multiplicative gate:
\begin{equation}
\text{score}_i = \text{sim}_i \times \bigl(1 + \sigma(\boldsymbol{\theta}^\top \mathbf{f}_i)\bigr),
\end{equation}
where $\text{sim}_i$ is cosine similarity, $\mathbf{f}_i$ is a vector of log-estimated persistent signals, and $\sigma$ is the sigmoid function. The gate ensures: (a) low-similarity nodes cannot be retrieved by importance alone, and (b) among similarly relevant nodes, persistent importance disambiguates.

\section{Related Work}

\textbf{Memory for LLM agents.}\quad Park et al.~\citep{park2023generative} introduced a three-term scoring function for generative agents. MemGPT~\citep{packer2023memgpt} manages memory via a virtual context hierarchy. Reflexion~\citep{shinn2023reflexion} uses self-reflection for episodic learning. None explicitly benchmark persistent-importance retrieval on long-horizon tasks.

\textbf{Retrieval-augmented generation.}\quad Standard RAG~\citep{lewis2020retrieval} uses dense similarity retrieval. Recent work adds reranking~\citep{glass2022re2g} and hybrid search. These focus on factual QA, not persistent-importance memory management.

\textbf{Attention and novelty.}\quad Predictive coding~\citep{friston2010free} and surprise-driven attention~\citep{itti2009bayesian} process prediction errors. Differential Transformer~\citep{ye2024diff} subtracts attention maps to reduce noise. Our work is complementary: MGA operates as an external retrieval layer, not a transformer modification.

\section{Method}

\subsection{Problem Setting}

A memory store contains $N$ nodes, each with content text, timestamp, and session ID. Given a query at time $t$, retrieve the top-$k$ nodes that answer the query. Gold labels exist only in an oracle and are never visible to the retrieval system.

\subsection{Signal Estimators}

All signals are estimated from the interaction log without access to gold labels:
\begin{itemize}[leftmargin=2em,itemsep=2pt]
\item \textbf{Recency}: $r_i = \exp(-(t_q - t_i) / \tau_r)$, with $\tau_r = 5$ days.
\item \textbf{Frequency}: $f_i = \log(1 + m_i) / \log(1 + m_{\max})$, where $m_i$ counts near-duplicate mentions (cosine similarity $> 0.75$).
\item \textbf{Unresolvedness}: binary heuristic detecting task-like nodes without subsequent closure matches.
\item \textbf{Goal relevance}: $g_i = \cos(\mathbf{e}_i, \mathbf{e}_{\text{goal}})$, normalized per retrieval call.
\item \textbf{Utility/Weight}: uninformative prior (0.5) in current benchmark.
\end{itemize}
All features are min-max normalized to $[0,1]$ per retrieval call.

\subsection{MGA Gate Retriever}

The persistent feature vector (excluding similarity) is:
\[
\mathbf{f}_i = [\text{recency}, \text{frequency}, \text{unresolved}, \text{utility}, \text{weight}, \text{goal\_rel}] \in [0,1]^6.
\]
The MGA gate score:
\begin{equation}
\text{score}_i = \text{sim}_i \times \bigl(1 + \sigma(\boldsymbol{\theta}^\top \mathbf{f}_i)\bigr).
\end{equation}
$\boldsymbol{\theta} \in \mathbb{R}^6$ is learned via logistic regression on training queries (gold = positive, sampled non-gold = negative). The gate value $\sigma(\boldsymbol{\theta}^\top \mathbf{f}_i) \in [0,1]$ means scores are bounded in $[\text{sim}_i, 2 \cdot \text{sim}_i]$.

\subsection{Baselines}

\begin{itemize}[leftmargin=2em,itemsep=2pt]
\item \textbf{Recency}: rank by $r_i$ only.
\item \textbf{Similarity}: rank by $\text{sim}_i$ only.
\item \textbf{GA}: $r_i + w_i + \text{sim}_i$ \citep{park2023generative}.
\item \textbf{MGA linear}: $\boldsymbol{\theta}^\top [\text{sim}, \mathbf{f}]$ (ablation without gate).
\end{itemize}

\section{Experimental Setup}

\subsection{Synthetic Benchmark}

We generate 10 synthetic worlds, each with: 25 sessions, 10--25 events per session ($\sim$450 nodes/world); 3 domains (research, software, design); 2--4 mid-log revisions (stale traps); 4--8 tasks with staggered open/close; 45\% chatter (noise); 5--8 templates per event type.

Query families: constraint recall, stale trap, open-loop, needle, noise resistance. Train/test: 40\%/60\% per world. $\boldsymbol{\theta}$ learned on train only.

\subsection{Embeddings and Evaluation}

Sentence-transformers (\texttt{all-MiniLM-L6-v2}) for similarity. GPU: NVIDIA T4 via Modal. Metrics: Recall@5, nDCG@5, Stale@5, Noise@5 on test queries.

\section{Results}

\subsection{Overall Performance}

\begin{table}[h]
\centering
\caption{Retrieval performance across 415 test queries (10 worlds). Best in \textbf{bold}.}
\vspace{0.5em}
\begin{tabular}{lcccc}
\toprule
\textbf{Retriever} & \textbf{Recall@5} & \textbf{nDCG@5} & \textbf{Stale@5} & \textbf{Noise@5} \\
\midrule
Recency & 0.029 & 0.085 & 0.000 & 0.364 \\
Similarity & 0.138 & 0.496 & 0.340 & 0.046 \\
GA & \textbf{0.209} & \textbf{0.697} & \textbf{0.000} & 0.070 \\
MGA gate & 0.174 & 0.562 & 0.265 & 0.029 \\
MGA linear & 0.195 & 0.614 & 0.205 & \textbf{0.017} \\
\bottomrule
\end{tabular}
\end{table}

GA achieves the highest overall Recall@5 and nDCG@5. Its stale advantage (0.000) is due to recency naturally filtering old nodes.

\subsection{Per-Family Breakdown}

\begin{table}[h]
\centering
\caption{Recall@5 by query family. \textbf{Bold} = best per row.}
\vspace{0.5em}
\begin{tabular}{lccccc}
\toprule
\textbf{Family} & \textbf{Recency} & \textbf{Sim} & \textbf{GA} & \textbf{MGA gate} & \textbf{MGA linear} \\
\midrule
Constraint & 0.036 & 0.151 & \textbf{0.257} & 0.191 & 0.220 \\
Needle & 0.000 & 0.143 & 0.000 & \textbf{0.143} & 0.000 \\
Open-loop & 0.000 & 0.097 & 0.000 & 0.111 & \textbf{0.222} \\
Noise resist & 0.000 & 0.050 & 0.000 & \textbf{0.075} & 0.025 \\
Stale trap & 0.036 & 0.135 & \textbf{0.275} & 0.187 & 0.230 \\
\bottomrule
\end{tabular}
\end{table}

Critical finding: \textbf{GA scores 0.000 on needle, open-loop, and noise-resistance families.} MGA achieves non-zero recall on all three.

\subsection{Learned Feature Importance}

\begin{table}[h]
\centering
\caption{Learned $\boldsymbol{\theta}$ (MGA gate), mean $\pm$ std across 10 worlds.}
\vspace{0.5em}
\begin{tabular}{lc}
\toprule
\textbf{Feature} & \textbf{Coefficient} \\
\midrule
Goal relevance & $2.92 \pm 0.87$ \\
Recency & $0.67 \pm 0.50$ \\
Frequency & $0.40 \pm 0.35$ \\
Unresolvedness & $0.21 \pm 0.29$ \\
Utility & $-0.01 \pm 0.01$ \\
Weight & $-0.01 \pm 0.01$ \\
\bottomrule
\end{tabular}
\end{table}

Goal relevance dominates ($\theta = 2.92$): knowing what the system is currently working on is the strongest persistent signal.

\subsection{Attention-Entropy Diagnostic}

We introduce an attention-entropy diagnostic to detect \emph{retriever lock}---when a retriever attends to a narrow set regardless of query content. We compute normalized Shannon entropy of the softmax-transformed score distribution:
\[
\hat{H} = \frac{-\sum_i p_i \log p_i}{\log N}, \quad p_i = \text{softmax}(s_i).
\]

\begin{table}[h]
\centering
\caption{Attention-entropy diagnostic (5 worlds, 29 test queries per retriever).}
\vspace{0.5em}
\begin{tabular}{lccc}
\toprule
\textbf{Retriever} & \textbf{Mean $\hat{H}$} & \textbf{Std $\hat{H}$} & \textbf{Top-5 Conc.} \\
\midrule
Recency & 0.992 & 0.001 & 0.050 \\
Similarity & 0.992 & 0.005 & 0.060 \\
GA & 0.983 & 0.005 & 0.078 \\
MGA gate & 0.976 & 0.014 & 0.093 \\
MGA linear & 0.650 & 0.131 & 0.467 \\
\bottomrule
\end{tabular}
\end{table}

No retriever shows pathological lock. MGA linear concentrates attention most (top-5 conc.\ = 0.467), explaining its targeted recall. MGA gate maintains high entropy (0.976) while still selectively boosting---the multiplicative structure amplifies without collapsing.

\section{Discussion}

\subsection{Why GA Fails on Persistent Tasks}

GA combines recency + importance + similarity additively. On needle queries, recency suppresses old nodes. On open-loop queries, the importance proxy is uninformative without explicit task-tracking. On noise queries, GA has no mechanism beyond similarity to distinguish signal from chatter.

MGA's gate addresses this: persistent features (especially goal-relevance and unresolvedness) amplify nodes that similarity alone identifies as marginally relevant.

\subsection{Limitations}

\begin{enumerate}[leftmargin=2em,itemsep=2pt]
\item \textbf{Synthetic data}: template-based logs. Real interactions have greater lexical diversity and ambiguity.
\item \textbf{Small training signal}: 4--5 training queries per world. More data would improve $\boldsymbol{\theta}$ estimation.
\item \textbf{Utility/weight uninformative}: insufficient training episodes for these features.
\item \textbf{Overall recall}: GA wins overall. MGA's advantage is specifically on long-horizon tasks.
\end{enumerate}

\subsection{Broader Significance}

MGA operationalizes a principle from existential systems theory: attend not only to what is \emph{similar} to the current query, but to what is \emph{persistently important} for ongoing operation. This parallels gap-based attention in biological systems, where survival-relevant mismatches receive priority processing regardless of surface similarity.

\section{Conclusion}

We introduced MGA, a retrieval mechanism that gates similarity by learned persistent importance:
\begin{enumerate}[leftmargin=2em,itemsep=2pt]
\item Standard approaches (similarity, GA) score zero on tasks requiring long-horizon persistent memory.
\item MGA achieves non-zero recall on all persistent-memory task families.
\item Goal-relevance is the dominant persistent signal ($\theta = 2.92$).
\item MGA has the lowest noise retrieval rate (0.017 vs GA's 0.070).
\item Attention-entropy diagnostic confirms healthy, non-locked behavior.
\end{enumerate}

MGA does not replace similarity-based retrieval. It fills the gap where similarity and recency fail: persistent, long-horizon, noise-resistant memory.

\begin{ack}
This work was conducted independently without external funding.
\end{ack}

\bibliographystyle{plainnat}
\bibliography{references}

\end{document}
